\input python_header

\begin{document}
%%%%%%%%%%%%%%%%%%%%%%%%%%%%%%%%%%%%%%%%%%%%%%%
\docheader{2026}{Python}{\S 2}{Martin Henz, Wang Yuyang, Lucas Ng, Jonas Ong}
%%%%%%%%%%%%%%%%%%%%%%%%%%%%%%%%%%%%%%%%%%%%%%%

\input python_intro

\section{Changes}

Compared to Python \S 1, Python \S 2 has the following changes:
\begin{itemize}
  \item \lstinline{None}: primitive linked list expression (empty linked list).
  \item Linked list library: Functions for creating, accessing and
    processing linked lists.
\end{itemize}

\input python_1_bnf

\input python_indentation
\input python_precedence
\input python_restrictions

\input python_single_assignment

\input python_names_lang
\input python_numbers

\input python_strings

\input python_typing_front
\input python_typing_middle_12
\input python_typing_back
\input python_typing_tail_12

The operators \texttt{\textbf{==}} and \texttt{\textbf{!=}} on pairs compare the corresponding
components; a component that is the same number as the other is equal to it, even if it is a
\texttt{\textbf{nan}}, as in Python:
\texttt{pair(math\_nan, 1) == pair(math\_nan, 1)} evaluates to \texttt{\textbf{True}},
whereas \texttt{math\_nan == math\_nan} evaluates to \texttt{\textbf{False}}.

\input python_standard

\input python_misc

\input python_math

\input python_linked_list_library

\input python_py_differences

\end{document}
